\chapter{Data, Code and Reproducibility}\label{appC-data-code}\label{app:data}

\section{Introduction}\label{sec:appC-intro}

This appendix has three roles: the register of every dataset the thesis consumes, with source, coverage, quality and licence; the description of every code module that computes a printed number; and the protocol by which a reader with the stated environment regenerates those numbers. Which windows are fitted, which months are reported and which policy is certified are pre-registered in Chapters~\ref{ch:usd} and~\ref{ch:zar}, not here, and the transition record of dates, spreads and conventions is kept in Chapter~\ref{ch:intro}.

Much of what follows is a specification rather than a record, so the status on 9~September 2026 is stated plainly. No market dataset has been downloaded and no terminal has been queried; every vendor identifier is either read from a primary document or marked to be confirmed. The code covers Chapters~\ref{ch:prelim} to~\ref{ch:sharpness} only: the multi-curve and time-changed SABR formulas, and the certificate with its worked example and figure.

\paragraph{Access classes.} One legend is used in both register tables. \emph{Class~F}: free public download, redistributable subject to the source's notice. \emph{Class~T}: available only through the university Bloomberg or LSEG terminal. \emph{Class~S}: available only from a sponsoring dealer under a data agreement. \emph{Class~C}: computed by the thesis code from rows of the other classes.

\paragraph{Licensing principles.} Four principles apply throughout and are not repeated at each dataset. Terminal data fall under the university's standard Bloomberg or LSEG agreement, which permits internal academic use but not redistribution, so no class~T file is shipped: the package carries the identifier, the export instruction and a checksum. Public data are shipped only where the source's terms allow it, as the New York Fed Terms of Use do; where the terms are unsettled, as for Reserve Bank exports and exchange files, the package ships a download script and a checksum. Quoting a few values from a public regulation or a published technical note is quotation of a document, not use of a licensed data service, and is how the fallback spreads of both currencies enter. Sponsor data never leave university-managed storage and enter only as the three aggregates of Section~\ref{sec:appC-vols}.

\section{USD data}\label{sec:appC-usd}% appC-usd-data.tex -- \input under \section{USD data} in appC-data-code.tex
% Every ticker below was read from a primary document (BISL fact sheet v7.2,
% ISDA, NY Fed, FRED, FCA, Federal Reserve Regulation ZZ, LSEG developer forum)
% or is marked TBC. Nothing here has been pulled from a terminal. Facts for which
% no primary document has been obtained carry \gap{Citation needed: ...}.

\subsection{Scope and notation}\label{appC:usd-purpose}

Chapter~\ref{ch:usd} needs, on every business day of its windows, (i) the USD LIBOR caplet and swaption smiles, (ii) the SOFR caplet and swaption smiles, (iii) the LIBOR--SOFR basis, (iv) the fallback spread $s$ by tenor, (v) SOFR discounting and forward curves, and (vi) the dates that bound the quoted markets. Tenors follow the ISDA convention O/N, 1W, 1M, 2M, 3M, 6M, 12M. Bloomberg tickers are written \texttt{TICKER <Index>} or \texttt{TICKER <Curncy>}; LSEG (Refinitiv) identifiers are written as RICs. Section~\ref{appC:usd-minimum} states the smallest set on which the chapter can run.

\subsection{Data constraints from the cessation calendar}\label{appC:usd-calendar}

The transition record (dates, spreads, conventions) is kept in Chapter~\ref{ch:intro}, and Chapter~\ref{ch:usd} fixes its fitting window $\Dfit$ and its deployment window from that calendar alone. The data constraints those choices must respect are three. A quoted USD LIBOR option surface exists to 31~December 2021, after which the market shrinks under the regulators' ``no new LIBOR'' guidance, and no LIBOR option market exists after the panel ended on 30~June 2023. A quoted SOFR option surface exists only from about November 2021 (the first month in which the vendor surface is marked quoted rather than derived is TBC, Section~\ref{appC:usd-sofr-vol}). The spread $s$ is fixed on 5~March 2021 and is constant thereafter.

\gap{Citation needed: FCA announcement of 5 March 2021 (cessation and non-representativeness of all LIBOR settings), ISDA statement of 5 March 2021 (Index Cessation Event, spread fixing date), FCA news of 3 July 2023 (panel end) and September 2024 (end of synthetic USD LIBOR), CFTC MRAC releases 8409-21 and 8449-21 (SOFR First, 8 November 2021 for non-linear products). None of these documents has been obtained and none has a bibliography entry on 9 September 2026, so the dates above stand as unverified until the primary texts are read and cited in Chapter~\ref{ch:intro}.}

\figplan{USD data availability, horizontal axis January 2018 to December 2024. Horizontal bars, one per series: USD LIBOR fixings (to 30 June 2023, hatched for synthetic LIBOR to 30 September 2024); LIBOR cap/floor and swaption surface (solid while quoted, hatched from 1 January 2022, absent after 30 June 2023); SOFR fixings (from 2 April 2018); SOFR cap/floor and swaption surface (hatched while vendor-derived, solid from the first quoted month, marked TBC); LIBOR--SOFR basis swaps (solid to mid-2022, hatched for the convergence period to 30 June 2023). Vertical lines at 5 March 2021 (spread fixing), 8 November 2021 (SOFR First, non-linear), 31 December 2021, 30 June 2023. The fitting and deployment windows chosen in Chapter~\ref{ch:usd} are shaded.}{Availability of the USD datasets of Table~\ref{tab:usd-register}. The figure tests one claim: that the windows pre-registered in Chapter~\ref{ch:usd} lie inside quoted, not derived, data on both benchmarks. Sources: the register table and the vendor flags recorded under the protocol of Section~\ref{sec:appC-protocol}.\label{fig:appC-usd-timeline}}

\subsection{Fallback spread adjustments}\label{appC:usd-spread}

For USD the constant $s$ of Definition~\ref{def:prelim-fallback} is the ISDA fallback spread adjustment: the five-year historical median of USD LIBOR minus SOFR compounded in arrears over the matching tenor, fixed on the Spread Adjustment Fixing Date of 5~March 2021 and published by Bloomberg Index Services Limited (BISL) as ISDA's calculation agent. Table~\ref{tab:usd-spread} records the values and the tickers that carry them.

\begin{table}[htbp]
\centering\small
\caption{USD LIBOR fallback spread adjustments fixed 5~March 2021 and the BISL tickers that carry them. The five numerical values are those enacted in Federal Reserve Regulation~ZZ (12~CFR~253.4(c)), which states that they equal the ISDA spread adjustments. Tickers from the BISL IBOR Fallbacks fact sheet v7.2 (26~September 2025), Appendix~2; all carry the \texttt{<Index>} key. TBC: to be read from the terminal.}
\label{tab:usd-spread}
\begin{tabular}{@{}lrlll@{}}
\toprule
Tenor & Spread (\%) & Spread ticker & Adjusted SOFR ticker & Fallback rate ticker \\
\midrule
O/N & 0.00644 & \texttt{YUS00ON} & \texttt{XSOFRON} & \texttt{VUS00ON} \\
1W & TBC & \texttt{YUS0001W} & \texttt{XSOFR1W} & \texttt{VUS0001W} \\
1M & 0.11448 & \texttt{YUS0001M} & \texttt{XSOFR1M} & \texttt{VUS0001M} \\
2M & TBC & \texttt{YUS0002M} & \texttt{XSOFR2M} & \texttt{VUS0002M} \\
3M & 0.26161 & \texttt{YUS0003M} & \texttt{XSOFR3M} & \texttt{VUS0003M} \\
6M & 0.42826 & \texttt{YUS0006M} & \texttt{XSOFR6M} & \texttt{VUS0006M} \\
12M & 0.71513 & \texttt{YUS001Y} & \texttt{XSOFR1Y} & \texttt{VUS001Y} \\
\bottomrule
\end{tabular}
\end{table}

The 3M value, $s=26.161$~bp, is the one the thesis uses for USD, because the caplet and swaption markets are written on 3M LIBOR. The fact sheet also gives \texttt{FUS0003M <Index>} for the all-in 3M fallback rate and names the underlying fixings \texttt{US0003M <Index>} (3M USD LIBOR) and \texttt{SOFRRATE <Index>} (SOFR). The \texttt{Y}-tickers are constant after 5~March 2021; the \texttt{X}-tickers give compounded-in-arrears SOFR over each LIBOR tenor with the two-day backward shift of the BISL Rule Book, which is the realised compounded rate $R(T,T+\Delta)$ of Definition~\ref{def:prelim-R} up to the day-count and shift conventions stated there. Quoting the five Regulation~ZZ values is use of a public regulation, not of Bloomberg's data; the package recomputes compounded SOFR from the New York Fed series (Section~\ref{appC:usd-ois}) rather than shipping any \texttt{X}-ticker history.

\gap{Citation needed: (a) Federal Reserve Board, Regulation~ZZ, 12~CFR Part~253, final rule (December 2022; Federal Register January 2023) for the five spread values and the statement that they equal the ISDA adjustments; the only Board document located for it, the press-release attachment of 19 July 2022, is the \emph{proposed} rule and not the final one, so the five values in Table~\ref{tab:usd-spread} rest on no obtained primary source and the Federal Register citation of the final rule must be read before they are cited; (b) BISL IBOR Fallbacks fact sheet v7.2, 26 September 2025, for every ticker in Table~\ref{tab:usd-spread}, the two-day shift and the licence terms (licence required for redistribution or usage of the adjusted rates, fee waived below USD~5bn total assets for a single rate set).}

\datanote{The 1W and 2M spread adjustments were fixed on the same date but are not in Regulation~ZZ (the LIBOR Act covers O/N, 1M, 3M, 6M and 12M only), and the BISL notice of 5~March 2021 that lists all seven (\url{https://assets.bbhub.io/professional/sites/10/IBOR_Fallbacks_LIBOR_Cessation_Announcement-2021-Mar.pdf}) rendered as images only on 9~September 2026. Read them from \texttt{YUS0001W} and \texttt{YUS0002M} at the terminal. Neither tenor is used in the thesis.}

\subsection{SOFR, USD LIBOR fixings and the SOFR OIS curve}\label{appC:usd-ois}

\emph{SOFR fixings} (class~F) are administered by the Federal Reserve Bank of New York and published each business day at about 08:00~ET, with a same-day revision at about 14:30~ET if the correction exceeds one basis point; the first publication was 3~April 2018 for value date 2~April 2018. The daily fields are the rate, the 1st, 25th, 75th and 99th percentiles, and volume. The 30-, 90- and 180-day compounded averages and the SOFR Index (initial value $1.00000000$ on 2~April 2018) are official from 2~March 2020 with indicative history to April 2018. Routes and mirrors are in Table~\ref{tab:usd-register}; the Terms of Use grant a licence to use, copy and distribute the content with an attribution notice, so the series may be shipped. Two data facts matter for daily scores: SOFR printed $5.25\%$ on 17~September 2019 against $2.43\%$ the day before, which the compounded 3M rate absorbs as roughly 3~bp, and the 08:00 and 14:30 prints can differ, so the final print is used.

\emph{USD LIBOR fixings} (class~T). No public history exists: FRED removed all IBA series on 31~January 2022 at IBA's request, and IBA's Historical LIBOR Rates page is view-only under a user agreement that forbids copying, redistribution and pricing use, with an Educational Usage Licence priced at USD~5{,}000 per benchmark per year. The thesis takes the fixings from the terminal and ships code that reads a user-supplied CSV.

\emph{SOFR OIS curve} (class~T). Chapter~\ref{ch:usd} needs SOFR discount factors and the forward compounded rate $F_t$ at each caplet start, bootstrapped from SOFR OIS par rates 1W to 50Y (annual fixed versus compounded SOFR, Act/360, T+2) and, at the short end, SOFR futures. No public daily history of SOFR OIS par rates exists; CME publishes one- and three-month SOFR futures settlements for view-only use, with history through a licensed DataMine order. For deployment dates the discount curve can be replaced by the realised SOFR path, because the certificate scores realised hedging error, not model value; only $F_t$ at fitting dates needs a curve.

\datanote{If the terminal export of the OIS strip is refused or incomplete, $F_t$ is bootstrapped from CME SOFR futures settlements out to two years (the longest caplet expiry used) with the convexity adjustment stated; the 2019--2023 settlement history must then be licensed through DataMine or transcribed at the terminal (\texttt{SFRA <Cmdty>} is the front contract, TBC).}

\gap{Citation needed: New York Fed SOFR publication page and Terms of Use (\url{https://www.newyorkfed.org/privacy/termsofuse}) for the publication times, revision threshold, percentile fields, Averages and Index dates, the licence clause and the 17 September 2019 print; FRED notice of 31 January 2022 on the removal of IBA series; ICE Historical LIBOR Rates User Agreement and the IBA fee schedule for the USD~5{,}000 figure; CME SOFR futures listing date 7 May 2018 and settlement terms. None has been obtained and none has a bibliography entry on 9 September 2026.}

\subsection{USD LIBOR cap/floor and swaption volatility surfaces}\label{appC:usd-libor-vol}

\emph{What is needed.} For each business day in $\Dfit$: (a) cap/floor normal volatilities on 3M USD LIBOR for expiries 1Y to 10Y at ATM and a fixed absolute grid (at least $1\%$, $1.5\%$, $2\%$, $2.5\%$, $3\%$, $4\%$, $5\%$, plus ATM$\pm25$, $\pm50$, $\pm100$~bp where quoted), from which caplet volatilities are stripped; (b) ATM swaption normal volatilities on the grid $\{$1M, 3M, 6M, 1Y, 2Y, 5Y$\}\times\{$1Y, 2Y, 5Y, 10Y$\}$ and, where available, the smile at $\pm25$, $\pm50$, $\pm100$~bp. Normal (basis-point) volatility is the market convention throughout; lognormal quotes, if downloaded, are converted with the forward of the same day, and the convention field is exported with every value.

\emph{Availability.} Class~T only: no public USD LIBOR cap or swaption volatility history exists, so this dataset is the one refusal that would stop the first half of Chapter~\ref{ch:usd}. Both vendors serve it as a composite volatility cube built from contributed broker quotes rather than as a quoted series, Bloomberg by ticker family and LSEG through volatility-surface endpoints rather than fixed RICs; the identifiers and their unresolved parts are in Table~\ref{tab:usd-register}. Daily closes are required from 2~January 2019 to 30~June 2023; Chapter~\ref{ch:usd} decides which part of that range is fitted and which is reported as a diagnostic of a dying market.

\emph{Known quality issues.} Composite surfaces are filled: \texttt{BVOL} and IPA interpolate and extrapolate where brokers did not quote, so a downloaded smile may be model output. The thesis records, for each day and point, whether the vendor flags the point as quoted or derived where the vendor exposes that flag; where it does not, strikes beyond $\pm100$~bp are treated as unquoted (protocol item~6, Section~\ref{sec:appC-protocol}). In 2020--2021 the ATM 3M forward was near zero and shifted-lognormal quotes may sit alongside normal quotes. Absolute strike grids are re-centred by vendors as rates move, so a strike present in 2019 may be absent in 2021; the caplet-stripping module, which is not written (Section~\ref{subsec:appC-certify}), will interpolate in strike only within the quoted range. After 31~December 2021 front-end LIBOR volatility is in the SOFR-first regime and LIBOR quotes are largely derived from SOFR quotes plus the fixed spread; this is a data fact that Chapter~\ref{ch:usd} takes into account when it fixes $\Dfit$.

\gap{Citation needed: the claim that off-ATM swaption quotes between August 2019 and April 2020 were largely missing from vendor surfaces was written as ``documented in the literature'' without naming a work; no source is on disk. Until a peer-reviewed or primary source is cited, the thesis relies only on the vendor quoted/derived flag for that period.}

\datanote{``The ATM 3M forward fell to about 0.2\%'' in 2020--2021 is a recollection of the LIBOR fixing level, unverified against the fixing series on 9 September 2026; the exact path is read from \texttt{US0003M} at the terminal.}

\subsection{SOFR cap/floor and swaption volatility surfaces}\label{appC:usd-sofr-vol}

\emph{What is needed.} The grids of Section~\ref{appC:usd-libor-vol} on compounded 3M SOFR (caps) and on SOFR OIS swaps (swaptions), daily, from the first date the vendor shows a quoted SOFR surface to 31~December 2024. The early part is the observed successor smile that Chapter~\ref{ch:usd} compares with the smile predicted from LIBOR under Chapter~\ref{ch:identifiability}; the later part is the deployment sample for the certificate.

\emph{Availability.} Class~T for the surfaces, on the same two cubes with the SOFR underlying selected and with tickers and RICs unresolved (Table~\ref{tab:usd-register}); one public source at the short end. The only public SOFR option market is CME options on three-month SOFR futures: daily settlements are viewable on cmegroup.com, history needs DataMine, and the implied smile on the first eight quarterly contracts is a genuine public successor smile out to two years, sufficient for the identifiability test at the 3M~$\times$~\{3M,\ldots,2Y\} expiries though not for the swaption grid. Until vendors had SOFR broker quotes the SOFR surface was the LIBOR surface shifted by $s$ with volatility carried across, so the first month in which the vendor marks the surface as quoted is the first item of the terminal checklist (Section~\ref{subsec:appC-checklist}); Chapter~\ref{ch:usd} fixes its reporting rule from that date.

\emph{Known quality issues.} Dealer caps trade on compounded-in-arrears SOFR while loan-market caps trade on CME Term SOFR, whose derivatives are restricted to end-user hedging; the thesis uses compounded caps only. The 2021 ISDA Definitions took effect for cleared trades on 2--3~October 2021 and quotes straddling that weekend may change day count or payment lag. Interdealer SOFR swaptions were physically settled from November 2021 pending a tradeable SOFR swap-rate benchmark; ISDA Supplement~88 (10~November 2021) and Settlement Matrix v3.0 set the cash-settlement rate to the USD SOFR ICE Swap Rate, so the settlement type is exported with each quote. SOFR forwards jump at FOMC dates; the continuous-forward SABR description of Chapter~\ref{ch:prelim} holds between meeting dates \citep{brace2024}, and Chapter~\ref{ch:usd} states its treatment of short-expiry caplets that span a meeting.

\datanote{Secondary sources, unverified: Clarus SDR analyses (10 November 2021 and the following week) reporting SOFR swaptions in SDR data from May 2021, USD~32bn against USD~80bn of LIBOR swaption notional on 8 November 2021 (13\%, 29\% the week after) and SOFR caps and floors at 19\% of cap/floor notional in the week of 8--12 November 2021; Risk.net (15 September 2022) reporting that SOFR swaption volatility data from major vendors remained unavailable and that banks mapped LIBOR surfaces to SOFR by shifting strikes by the spread; a Deriscope example of 30 December 2022 for the layout of the SOFR cap table. None is a primary source and no fact in the thesis rests on them; they motivate the TBC checklist only.}

\gap{Citation needed: ISDA Supplement 88 (10 November 2021) and Settlement Matrix v3.0; ISDA guidance on the USD LIBOR ICE Swap Rate (28 November 2022, updated 13 April 2023); the 2021 ISDA Definitions effective date for cleared trades; CME Term SOFR launch (21 April 2021) and ARRC recommendation (29 July 2021); CME SOFR options listing date. None has been obtained and none has a bibliography entry. Separately, \citet{brace2024} is cited above only for the qualitative feature that SOFR forwards jump at scheduled policy dates and are diffusive between them; the published article was not obtained on 9 September 2026, so it carries no locator and no quantitative claim of this thesis rests on it.}

\subsection{LIBOR--SOFR basis swaps}\label{appC:usd-basis}

\emph{What is needed.} Daily par spreads of 3M USD LIBOR versus compounded SOFR basis swaps, 1Y to 30Y, 2019 to April 2023, from which Chapter~\ref{ch:usd} estimates the level, the volatility $\eta$ and the correlations $\rho_{BF}$, $\rho_{B\alpha}$ of the basis $B_t$ under the physical measure; Chapter~\ref{ch:identifiability} shows the option market does not identify them, and the historical estimate is the ambiguity set the certificate takes as input.

\emph{Availability.} Class~T for quoted spreads, on an LSEG contributed curve whose history depth is unresolved and on Bloomberg basis-swap tickers whose roots are unresolved (Table~\ref{tab:usd-register}). Two substitutes exist: (i) the realised basis, 3M LIBOR fixing minus compounded SOFR over the same period minus $s$, computable from Section~\ref{appC:usd-ois} and redistributable in the part that rests on the New York Fed series; (ii) the Fed-funds--SOFR and LIBOR--Fed-funds pairs on the same terminals. The realised basis is what $B_t$ denotes; the quoted basis swap is its risk-neutral expectation and is needed only as a check on drift.

\emph{Known quality issue.} Quoted basis spreads compress mechanically toward $s$ as June 2023 approaches, because CME and LCH converted cleared LIBOR swaps at the fixed spread in April--May 2023; the last twelve months of the quoted series are a convergence trade, not a free market. Chapter~\ref{ch:usd} fixes the estimation window for $\eta$ with this constraint; the realised basis is available for the whole of 2018--2023.

\gap{Citation needed: CME (21 April 2023) and LCH (22 April and 20 May 2023) conversion notices for cleared USD LIBOR swaps; LSEG developer-forum threads for the RICs \texttt{USD3MFSR=}, \texttt{USDSOFR=}, \texttt{USDSROIS=} and \texttt{USDSR3LBS=TPSR} (forum posts are not primary documents, so these RICs remain TBC until confirmed in Workspace).}

\subsection{Dataset register}\label{appC:usd-register}

\begin{table}[htbp]
\centering\scriptsize
\caption{USD datasets for Chapter~\ref{ch:usd}. Access classes F, T, S, C as defined in Section~\ref{sec:appC-intro}; TBC means to be confirmed at the terminal. ``In package'' follows the protocol of Section~\ref{sec:appC-protocol}: yes, hash-only, or script (the package ships the download script and the hash, not the file).}
\label{tab:usd-register}
\begin{tabular}{@{}p{2.4cm}p{3.0cm}p{2.3cm}p{2.1cm}cp{1.6cm}@{}}
\toprule
Dataset & Bloomberg; LSEG & Public source & Range (daily unless stated) & Class & In package \\
\midrule
Fallback spread $s$ by tenor & \texttt{YUS0003M} etc.\ \texttt{<Index>} & Reg.~ZZ 12~CFR~253.4(c) (five tenors) & fixed 5~Mar~2021 & F & yes (five values) \\
SOFR fixings & \texttt{SOFRRATE <Index>}; \texttt{USDSOFR=} (TBC) & NY Fed API; FRED \texttt{SOFR} & 2018-04-02-- & F & yes, with notice \\
SOFR averages, index, percentiles, volume & --- & NY Fed API (\texttt{/api/rates/secured/sofr/}); FRED \texttt{SOFR30/90/180DAYAVG}, \texttt{SOFRINDEX}, \texttt{SOFR75}, \texttt{SOFRVOL} & 2018-04-02-- (official from 2020-03-02) & F & yes, with notice \\
USD LIBOR fixings & \texttt{US0003M <Index>}; \texttt{USD3MFSR=} (TBC) & none (FRED removed 31~Jan~2022; IBA page view-only) & to 2023-06-30 (synthetic to 2024-09-30) & T & hash-only \\
Compounded SOFR by LIBOR tenor & \texttt{XSOFR3M <Index>} & recomputed from NY Fed SOFR & 2018-- & C & yes (code) \\
SOFR OIS par rates & \texttt{USOSFR\ldots\ <Curncy>} (TBC; \texttt{ICVS} curve number TBC); \texttt{USDSROIS=} (TBC) & none & 2019-- & T & hash-only \\
SOFR futures settlements & \texttt{SFRA <Cmdty>} (TBC) & CME settlements page (view only); DataMine (licensed) & 2018-05-07-- & T & hash-only \\
LIBOR cap/floor vols & \texttt{VCUB}/\texttt{BVOL} or \texttt{BBIR}, tickers TBC, history via \texttt{HP}/\texttt{BDH}; LSEG IPA \texttt{NormalizedVolatility} & none & 2019-01-02 to 2023-06-30 & T & hash-only \\
LIBOR swaption vols & \texttt{VCUB}; \texttt{USSN\ldots\ <Curncy>} (TBC); IPA & none & as above & T & hash-only \\
SOFR cap/floor vols & \texttt{VCUB}/\texttt{BVOL}, SOFR underlying, tickers TBC; IPA & CME SOFR futures options (short end) & first quoted month (TBC, $\ge$~Nov~2021) to 2024-12-31 & T & hash-only \\
SOFR swaption vols & \texttt{VCUB}, tickers TBC; IPA & none & as above & T & hash-only \\
LIBOR--SOFR basis swaps & \texttt{SWPM}/\texttt{ICVS} roots TBC; \texttt{USDSR3LBS=TPSR}, points \texttt{USDSR3L\{3M\ldots50Y\}=TPSR} (depth TBC) & realised basis from fixings & 2019--2023 quoted; 2018--2023 realised & T / C & hash-only / yes (code) \\
\bottomrule
\end{tabular}
\end{table}

\subsection{Minimum viable data}\label{appC:usd-minimum}

Chapter~\ref{ch:usd} can be written on three terminal exports and two public downloads. The terminal exports are: the daily 3M USD LIBOR fixing from 2~January 2018 to 30~June 2023; the daily \texttt{VCUB} caplet-volatility strip on 3M LIBOR at ATM and $\pm50$, $\pm100$~bp for expiries 3M to 2Y from 2~January 2019 to 31~December 2021; and the same strip on compounded SOFR from the first quoted month to 31~December 2024. The public downloads are the New York Fed SOFR series from 2~April 2018 and the five Regulation~ZZ spread adjustments. From these five inputs the realised compounded rate $R(T,T+\Delta)$, the realised basis $B_t$, the constant $s=26.161$~bp, the fitted SABR smiles, the predicted SOFR smile with its $\pm\eta$ interval and the realised hedging-error scores of the certificate all follow. Everything else in Table~\ref{tab:usd-register} improves the chapter: the swaption grid extends the identifiability test beyond two years, the OIS strip replaces the futures-implied forward, the quoted basis swaps check the drift of $B_t$, and the CME options give a public cross-check of the SOFR smile. If the terminal export of volatilities is refused, the chapter cannot be written on public data alone, because no public LIBOR option surface exists; the CME SOFR options then support only the second half of the chapter, and Chapter~\ref{ch:intro} must say so.

\datanote{Everything in this section is a specification written on 9~September 2026 from public documents. No terminal has been queried. Every TBC is resolved in the single terminal session of Section~\ref{subsec:appC-checklist}, and no result of Chapter~\ref{ch:usd} rests on an identifier still marked TBC here.}

\section{ZAR data}\label{sec:appC-zar}% appC-zar-data.tex --- \input under \section{ZAR data} in appC-data-code.tex
% Every URL was opened on 9 September 2026 unless the surrounding \datanote says
% otherwise. Access classes F/T/S/C are defined in the appendix introduction.

This section lists every South African dataset that Chapter~\ref{ch:zar} consumes, with its access class, route, format and coverage. Access dates are 9~September 2026. Table~\ref{tab:appC-zar} summarises; Chapter~\ref{ch:zar} states which data regime (sponsor or public) it was written under.

\subsection{ZARONIA daily fixings}\label{sec:appC-zaronia}

\emph{Class~F.} The South African Reserve Bank (SARB) administers ZARONIA and publishes the fixing for each South African business day at 10:00, with a republication window to 12:00 for corrections. Rates for 1~November 2022 to 3~November 2023 are labelled ``observation period rates'', rates from 4~November 2023 ``official publication rates'', and rates dated before 31~October 2022 ``pre-publication or proxy rates''. All three classes are served from one page,
\begin{quote}\footnotesize\raggedright\url{https://www.resbank.co.za/en/home/what-we-do/financial-markets/south-african-overnight-index-average}
\end{quote}
through a date-range selector that exports CSV, XML or PDF; there is no static file URL, so the protocol of Section~\ref{sec:appC-protocol} records the selected range, the earliest date returned and the download timestamp with every export. The SARB also publishes compounded ZARONIA period averages (1~week, 1, 3, 6, 9 and 12 months) and a compounded index with base $100$ on 1~November 2022, in the same formats, at
\begin{quote}\footnotesize\raggedright\url{https://www.resbank.co.za/en/home/what-we-do/financial-markets/compounded-zaronia-period-averages-and-index}
\end{quote}
The thesis computes $R(T,T+\Delta)$ of Definition~\ref{def:prelim-R} from the daily fixings and uses the period averages only as a check on day-count and business-day handling. Proxy data are available from 4~January 2016 \citep[\S5.1, p.~9]{sarb2025fallback}; the SAFEX overnight rate serves as the proxy for 2000 to 2015 \citep[\S2, p.~4]{sarb2025ois}.

\gap{Citation needed: the SARB ZARONIA benchmark page and statement of methodology for the 10:00 publication time, the 12:00 republication window, the publication start date and the three history classes; the page was read on 9 September 2026 but has no bibliography entry, and the statement of methodology itself was not obtained.}

\datanote{Two publication start dates are in circulation and must be reconciled at download: the SARB page labels rates from 1~November 2022 as the observation period, while \citet[\S5.3.1, p.~11]{sarb2025fallback} write that ZARONIA ``has been published by the SARB since 1 August 2022''. Whether the full proxy series from January 2016 is exportable from the selector was not verified on 9~September 2026; the protocol records the earliest date the selector returns. If proxy fixings before November 2022 are needed and are not exportable, they are requested from \texttt{sarb-benchmarks@resbank.co.za}; failing that, the 2016--2022 proxy enters only through the synthetic curve of Section~\ref{sec:appC-synthetic}, which embeds it.}

\subsection{SARB synthetic historical ZARONIA OIS curve}\label{sec:appC-synthetic}

\emph{Class~F.} The SARB Market Practitioners Group's Derivatives Workstream published \emph{Historical estimation of the ZARONIA OIS curve} (document dated 6~October 2025, posted 9~October 2025) \citep[\S1, p.~4]{sarb2025ois} with five Excel binary workbooks (\texttt{.xlsb}): \texttt{Benchmark\_Rates}; \texttt{Forecast\_Curve\_Expiry\_Dates} and \texttt{Forecast\_Curve\_\allowbreak Discount\_Factors}; \texttt{Discount\_Curve\_Payment\_Dates} and \texttt{Discount\_Curve\_\allowbreak Discount\_Factors}. Landing page:
\begin{quote}\footnotesize\raggedright\url{https://www.resbank.co.za/en/home/publications/publication-detail-pages/Financial-Markets/Committees/MPG/MPG-Related-pages/2025/historical-estimation-of-the-zaronia-ois-curve}
\end{quote}
Construction, as the paper states: the input is the JSE's daily bootstrapped nominal swap zero curve (deposits, FRAs and swaps on 3-month JIBAR), available from early 2005 to the end of 2024 \citep[\S2, p.~4]{sarb2025ois} with four missing dates, 24~Dec~2007, 2~Jan~2009, 24~Nov~2011 and 21~May~2013 \citep[\S2, fn.~2, p.~4]{sarb2025ois}. Daily 3-month JIBAR forwards are read off that curve and shifted by a tenor-dependent spread equal to the trailing median of the spot 3-month JIBAR minus ZARONIA spread over a window matching the forward's tenor \citep[\S3, p.~6]{sarb2025ois}; the shifted forwards are compounded into discount factors, and par rates are produced for spot-starting OIS from 1 to 12 months in monthly steps and from 15 months to 30 years in 3-month steps. The Workstream built the dataset for VaR and expected-shortfall scenario history \citep[\S1, p.~4]{sarb2025ois}. Because its JIBAR--ZARONIA basis is a deterministic rolling median, the curve must never be used to estimate $\eta$, $\rho_{BF}$ or $\rho_{B\alpha}$; Chapter~\ref{ch:zar} states what it is used for.

\subsection{JIBAR fixings, 1M/3M/6M/12M}\label{sec:appC-jibar}

\emph{Class~F.} SARB is the JIBAR administrator; its JIBAR page exports the fixing history as CSV, XML or PDF through the same selector mechanism as ZARONIA:
\begin{quote}\footnotesize\raggedright\url{https://www.resbank.co.za/en/home/what-we-do/statistics/johannesburg-interbank-average-rate}
\end{quote}
The JSE client portal keeps the longer archive under \emph{Downloadable Files~$\to$ Safex~$\to$ mtmdata}: a current \texttt{jbarRates.xls}, one folder per year from 2010 to 2026, and a ``Jibar Archive'' folder:
\begin{quote}\footnotesize\raggedright\url{https://clientportal.jse.co.za/downloadable-files?RequestNode=%2Fsafex%2Fmtmdata}
\end{quote}
JIBAR is published in 1-, 3-, 6-, 9- and 12-month tenors; all five cease permanently immediately after their final publication on 31~December 2026 \citep[p.~1]{isda2025}, so the fixing series is closed at that date and the JSE archive is the source of record for the 2000--2026 history that the five-year spread window and the Workstream's spread history draw on.

\datanote{The earliest date served by the SARB selector and the exact tenor columns in the JSE yearly files were not inspected on 9~September 2026 (the JSE portal was reached only through its search index); both are recorded at download. Neither the SARB JIBAR page nor the JSE portal has a bibliography entry.}

\subsection{Fallback credit adjustment spreads}\label{sec:appC-cas}

\emph{Class~F for the values; T for the daily fallback series.} Bloomberg Index Services Limited (BISL) is the ISDA fallback calculation agent for JIBAR. The SARB announcement of 3~December 2025 is the Index Cessation Event and the Spread Adjustment Fixing Date for all tenors \citep[Sect.~1.1]{isda2025}; BISL's technical note of 22~December 2025 \citep[p.~1]{bisl2025} fixes the spreads in Table~\ref{tab:appC-cas}. The methodology is the five-year historical median of JIBAR minus compounded ZARONIA recommended by the MPG \citep[\S1, p.~3]{sarb2025fallback}, with the lookback window specified relative to the fixing date after subtracting the tenor and the two-banking-day payment delay \citep[\S4.1, p.~7]{sarb2025fallback}.

\begin{table}[htbp]\centering\small
\caption{Fixed JIBAR fallback spread adjustments, Spread Adjustment Fixing Date 3~December 2025, from the BISL technical note of 22~December 2025 \citep[p.~1, Table]{bisl2025}.}\label{tab:appC-cas}
\begin{tabular}{lllr}
\toprule
Tenor & Bloomberg ticker & Spread (\%) & Spread (bp)\\
\midrule
1 month  & \texttt{YJIBA1M Index} & 0.1142 & 11.42\\
3 months & \texttt{YJIBA3M Index} & 0.1619 & 16.19\\
6 months & \texttt{YJIBA6M Index} & 0.4323 & 43.23\\
9 months & \texttt{YJIBA9M Index} & 0.5830 & 58.30\\
12 months& \texttt{YJIBA1Y Index} & 0.7410 & 74.10\\
\bottomrule
\end{tabular}
\end{table}

The thesis needs only the constant $s$ of Definition~\ref{def:prelim-fallback}, the 3-month value $0.1619\%$ \citep[p.~1, Table]{bisl2025}; the other tenors enter only Appendix~\ref{appA-tenor}. The daily all-in fallback rate, adjusted ZARONIA and spread series that BISL publishes at the terminal (\texttt{FBAK <GO>}) are licensed and not required: the thesis recomputes compounded ZARONIA from the free fixings and adds the fixed spread. Quoting the five values is quotation of a public technical note, not use of Bloomberg's licensed series.

\datanote{\citet{bisl2025} is a two-page technical note, read on 9~September 2026 at \url{https://assets.bbhub.io/professional/sites/27/IBOR-Fallbacks-ZAR-JIBAR_Cessation_Technical-Note_251222.pdf} and marked ``for information purposes only''; the five values and the five \texttt{YJIBA} tickers are on its single table page, which is the locator cited above. It does not state the five-year median methodology, which is why \citet{sarb2025fallback} is cited for that and not this note.}

\subsection{ZARONIA OIS and JIBAR swap curves}\label{sec:appC-curves}

\emph{Class~F to end-2024; T from 2025.} For dates up to 31~December 2024 the JIBAR swap curve is the JSE nominal swap zero curve (daily, in the same \texttt{mtmdata} tree as the JIBAR files) and the ZARONIA OIS curve is the Workstream's synthetic curve of Section~\ref{sec:appC-synthetic}. From January 2025 the live ZARONIA OIS market is quoted at the terminals, and the thesis needs end-of-day par OIS rates at the Workstream's pillars (1M--12M monthly, 15M--30Y quarterly) together with the JIBAR swap curve until JIBAR pricing screens close. The two curves give $F_t$ and $L_t$ for every accrual window and the discount factors for every payment date.

\datanote{Unverified on 9~September 2026: that ZARONIA OIS have been cleared at LCH since September 2024 and listed on Bloomberg SEF from 13~March 2025 (no primary source on disk); Bloomberg curve identifiers and LSEG RICs for the ZARONIA OIS and JIBAR swap curves (not found in any public document; TBC at the terminal); whether the JSE will publish a ZARONIA-based zero curve, and from what date. The BISL tickers of Table~\ref{tab:appC-cas} are the only ZAR identifiers verified from a primary source. A commercial provider (BlueGamma) advertises a downloadable ZARONIA forward curve; it is not a primary source and the thesis does not rely on it.}

\subsection{JIBAR cap/floor and swaption volatility surfaces}\label{sec:appC-vols}

\emph{Class~S.} No public or terminal source of ZAR cap, floor or swaption volatilities exists that the thesis can cite: the JIBAR cap market is quoted by dealers, not on a screen, and no evidence of quoted ZARONIA cap, floor or swaption screens was found as of 9~September 2026 (absence of evidence, not proof of absence; re-checked before submission). The MPG conventions paper says these products ``will in future be quoted on screens and/or via interbank broking agents'' \citep[\S3, p.~5]{sarb2024conv}.

What a sponsoring dealer would provide, under a data agreement that restricts publication to aggregates:
\begin{enumerate}[nosep]
\item End-of-day JIBAR cap/floor implied volatilities, normal (basis-point) quotation, for each business day from an agreed start date to 31~December 2026: one CSV per date with columns \emph{expiry}, \emph{strike} (absolute or offset from the at-the-money forward), \emph{normal volatility}, \emph{quote type} (mid, or bid and ask), \emph{source} (desk mark or consensus).
\item The same for ZARONIA caps and floors from the first date the desk marks them; the MPG's recommended strike grid is ATM and ATM~$\pm25$, $\pm50$, $\pm75$, $\pm100$, $\pm150$, $\pm200$~bp \citep[\S7.6, p.~25]{sarb2024conv}.
\item A swaption cube (expiry $\times$ tail $\times$ strike offset) in the same layout, JIBAR-referenced until cessation and ZARONIA-referenced thereafter.
\item The desk's own strike and expiry grids and interpolation rule, so that the thesis reproduces the desk's surface rather than resampling it.
\end{enumerate}
The thesis publishes no raw quote. Chapter~\ref{ch:zar} reports only (i) fitted SABR triples $(\alpha_0,\rho,\nu)$ per expiry, (ii) at-the-money volatility and skew differences between the JIBAR and ZARONIA surfaces in basis points, and (iii) monthly averages of these; none recovers a dealer's quote. The agreement should permit exactly these three aggregates and be signed by month~12 of the programme (the programme timetable places the sponsor decision at months 10--12).

\paragraph{Governance of sponsor data.} Raw sponsor files are held on university-managed encrypted storage under the candidate's account, with access limited to the candidate and the supervisors named in the agreement; no copy is placed in the reproducibility package or in any cloud service outside the university's. Retention is for the duration of the candidature plus the examination and revision period, after which the files are destroyed and the destruction recorded, unless the agreement fixes an earlier date. The agreement provides that the thesis examiners may inspect the raw quotes under the same confidentiality terms, so that a restriction on publication is not a restriction on examination. The data contain no personal information, so the Protection of Personal Information Act is not engaged. The three permitted aggregates above are the only outputs that leave the storage.

\gap{The university research-data management and intellectual-property policies under which the sponsor agreement is signed are to be named here with their current version and date; not yet identified on 9~September 2026.}

\subsection{Converted-book hedging P\&L from public fixings}\label{sec:appC-pnl}

\emph{Class~F, except one input.} The block score $S$ of Definition~\ref{def:score} for the South African book is the realised hedging error of a fixed policy $\pi$ over a block of monthly rebalancing periods; Chapter~\ref{ch:zar} defines the policy, the block length and the fallback when no sponsor signs. The inputs are:
\begin{enumerate}[nosep]
\item the realised compounded rate $R(T,T+\Delta)$ from the daily ZARONIA fixings (Section~\ref{sec:appC-zaronia});
\item the JIBAR fixing $J_T$ for windows that fixed before cessation (Section~\ref{sec:appC-jibar});
\item the constant $s=0.1619\%$ (Section~\ref{sec:appC-cas});
\item the forwards $F_t$, $L_t$ and discount factors $P(t,T+\Delta)$ from the curves of Section~\ref{sec:appC-curves} at each monthly rebalancing date;
\item one volatility level per expiry, frozen at the last date a JIBAR surface is available (Section~\ref{sec:appC-vols}; the only input of class~S).
\end{enumerate}
The record starts at the first monthly date after the LCH switch of ZAR discounting to ZARONIA on 11~April 2026 and runs to the last date before submission.

\datanote{Record length and feasibility. At most about twenty-four monthly hedging-error observations are available by month~30 of the programme; the count assumes month~0 is the proposal approval and a record starting in May 2026, and it is arithmetic, not data. With $N$ monthly periods and block length $b$ there are $n=\lfloor N/(2b)\rfloor$ calibration blocks, and the threshold of Definition~\ref{def:prelim-split} is finite only if $k=\lceil(1-\alpha)(n+1)\rceil\le n$, that is $n\ge(1-\alpha)/\alpha$, so $n\ge9$ at $\alpha=0.1$. Hence $N\ge18$ at $b=1$ and $N\ge36$ at $b=2$: on a record of about twenty-four months only $b=1$ yields a certificate, and the mixing class assumed in Chapter~\ref{ch:zar} must then justify a one-month gap. The LCH date rests on member updates of 17~December 2025 and 26~February 2026 that have not been obtained, so it has no bibliography entry; it is cited in Chapter~\ref{ch:intro} once the primary notice is read.}

\subsection{Summary}\label{sec:appC-zar-summary}

\figplan{ZAR data availability, horizontal axis January 2005 to June 2027. Horizontal bars: ZARONIA fixings in three shades (proxy to 31 October 2022, observation 1 November 2022 to 3 November 2023, official from 4 November 2023; proxy start marked TBC); JIBAR fixings to 31 December 2026; JSE nominal swap zero curve from 3 January 2005; SARB synthetic ZARONIA OIS curve 3 January 2005 to 31 December 2024; live ZARONIA OIS curve from 2025 (start TBC); sponsor volatility window (agreed start to 31 December 2026, dashed); monthly converted-book score record from May 2026 (dashed to submission). Vertical lines at 3 December 2025 (spread fixing), 11 April 2026 (LCH discounting switch), 31 December 2026 (JIBAR cessation).}{Availability of the ZAR datasets of Table~\ref{tab:appC-zar}. The figure tests the claim that the monthly score record of Section~\ref{sec:appC-pnl} lies entirely on official ZARONIA fixings and live or synthetic curves, and shows how short that record is relative to the feasibility arithmetic recorded there. Sources: the register table and the SARB, JSE, BISL and LCH documents cited in this section.\label{fig:appC-zar-timeline}}

\begin{table}[htbp]\centering\small
\caption{ZAR datasets for Chapter~\ref{ch:zar}: access class, source, coverage, licence and package status as recorded on 9~September 2026; rows marked TBC are unverified. Classes F, T, S, C as in Section~\ref{sec:appC-intro}.}\label{tab:appC-zar}
\scriptsize
\begin{tabular}{@{}p{2.5cm}cp{2.2cm}p{2.9cm}p{2.5cm}p{1.6cm}@{}}
\toprule
Dataset & Class & Source & Coverage & Licence & In package\\
\midrule
ZARONIA daily fixings & F & SARB benchmark page, selector (CSV/XML/PDF) & Official from 4 Nov 2023; observation from 1 Nov 2022; proxy earlier (extent TBC) & SARB website terms (redistribution TBC) & script + hash\\
Compounded ZARONIA averages and index & F & SARB, selector & 1W--12M from 1 Nov 2022 & as above & script + hash\\
Synthetic ZARONIA OIS curve & F & SARB DWS, Oct 2025, five \texttt{.xlsb} & Daily 3 Jan 2005 -- 31 Dec 2024, four dates missing; pillars 1M--12M, 15M--30Y & SARB copyright; link, not copy & hash-only\\
JIBAR fixings 1M--12M & F & SARB selector; JSE \texttt{mtmdata} (\texttt{.xls}) & JSE yearly files 2010--2026 plus archive; ends 31 Dec 2026 & SARB terms; JSE market-data terms (TBC) & script + hash\\
Fallback spreads by tenor & F & BISL technical note 22 Dec 2025 & Fixed 3 Dec 2025 & public note; five values quoted & yes\\
Daily BISL fallback rates & T & Bloomberg \texttt{FBAK}, \texttt{YJIBA*} & not required & terminal agreement & no\\
JIBAR swap zero curve & F & JSE \texttt{mtmdata}, daily files & 2005 -- present & JSE market-data terms (TBC) & hash-only\\
ZARONIA OIS par curve, live & T & Bloomberg/LSEG (IDs TBC) & 2025 -- present & terminal agreement & hash-only\\
JIBAR and ZARONIA cap/floor and swaption vols & S & sponsoring dealer, per-date CSV & agreed start to 31 Dec 2026 & data agreement; aggregates only & no\\
Converted-book hedging P\&L & C & thesis code from the rows above & monthly from May 2026 & --- & yes\\
\bottomrule
\end{tabular}
\end{table}

\gap{Citation needed: the SARB website terms of use and the JSE client-portal market-data terms, to settle whether SARB exports and JSE files may be redistributed; until settled the package ships the download script, the selector settings and the SHA-256 hash for those rows, not the files.}

If no sponsor signs by month~12, Chapter~\ref{ch:zar} is written from the class~F rows of Table~\ref{tab:appC-zar} and labels its reserve decomposition accordingly; the decision and its consequences are stated there, not here.

\section{Code architecture}\label{sec:appC-code}\subsection{The \texttt{rates} package}\label{subsec:appC-rates}

The multi-curve bootstrap and the time-changed SABR caplet pricer of Chapter~\ref{ch:prelim} are implemented as a small Python package, \texttt{thesis/code/rates/}, with a test file \texttt{thesis/code/tests/} \texttt{test\_rates.py}. The package depends on \texttt{numpy} only; the normal distribution function is taken from the standard library (\texttt{math.erf}), so no \texttt{scipy} installation is needed. The code is a reference implementation of the formulas of Chapter~\ref{ch:prelim}, not a production library: it has three modules and no calibration routine.

\paragraph{\texttt{curve.py}: discount curve and forward compounded rate.}
\texttt{bootstrap\_ois} returns discount factors $P(0,t_i)$ at increasing pillar maturities $t_1<\dots<t_n$ from the par rates $R_i$ of overnight index swaps maturing at $t_i$. The floating leg of an OIS is worth $1-P(0,t_i)$ at inception; the fixed leg is taken to pay $R_i(t_j-t_{j-1})$ at every pillar $t_j\le t_i$, so that the bootstrap is the closed recursion
\[
P(0,t_i)=\frac{1-R_i\sum_{j<i}(t_j-t_{j-1})P(0,t_j)}{1+R_i\,(t_i-t_{i-1})},
\]
with no root finding. The convention that fixed coupons fall on the pillar grid is exact when the pillars beyond one year are annual, which is the case for the curves used in Chapter~\ref{ch:usd}; a general coupon schedule would replace the annuity sum and nothing else. \texttt{discount} interpolates $\log P(0,t)$ linearly in $t$ between pillars, that is, with a constant continuously compounded forward rate on each interval, with $P(0,0)=1$ at the short end and a flat forward beyond the last pillar. \texttt{forward\_compounded} returns $F_0=(P(0,T)/P(0,T+\Delta)-1)/\Delta$, the identity $1+\Delta F_0=P(0,T)/P(0,T+\Delta)$ of Section~\ref{sec:prelim-compounded} under the continuous approximation.

\paragraph{\texttt{clock.py}: decay profile and clock.}
\texttt{g\_linear} is the profile of Assumption~\ref{ass:prelim-linear}, $g(t)=1$ for $t\le T$ and $(T+\Delta-t)/\Delta$ on $(T,T+\Delta]$. \texttt{tau} computes the clock $\tau(t)=\int_0^tg(u)^2\,du$ by composite Simpson quadrature on $[T,t]$, adding $T$ for the part where $g=1$. The value $\tau(T+\Delta)=T+\Delta/3$ is deliberately not hard-coded: it is a consequence of the linear-decay assumption, and the test verifies it numerically to $10^{-12}$ (Simpson's rule is exact for the quadratic $g^2$, so the agreement is to rounding). Any other profile can be passed as \texttt{g}; the pricer then uses its clock, and the factor $\Delta/3$ is replaced by whatever $c_g$ the profile implies, as in Example~\ref{ex:prelim-cg}.

\paragraph{\texttt{sabr.py}: Hagan volatility, Bachelier price, two caplet variants.}
\texttt{hagan\_normal\_vol} is formula \eqref{eq:prelim-hagan-N}, reducing to \eqref{eq:prelim-hagan-N0} at $\beta=0$; at $\zeta=0$ it uses $\zeta/x(\zeta)=1$. \texttt{bachelier\_call} is the undiscounted Bachelier call of Definition~\ref{def:prelim-implied}. \texttt{caplet\_time\_changed} is Lemma~\ref{lem:prelim-tc} in code: the Hagan volatility at the clock $\tau(T+\Delta)$, the Bachelier price at that clock, times $P(0,T+\Delta)$, per unit accrual. Both the compounded and the JIBAR caplet are priced under $\Q^{T+\Delta}$, and the code contains no convexity adjustment. \texttt{caplet\_nondecaying\_volvol} is the variant of Proposition~\ref{prop:prelim-novol} and is labelled in the source as an approximation, because no closed form exists for time-dependent vol-of-vol. Its two accountings are the two of the gap box after Proposition~\ref{prop:prelim-nueff}: \texttt{accounting="total"} uses the total-variance identification $\nu_{\mathrm{eff}}^2=\nu^2(T+\Delta)/\tau^*$ with $\rho$ unchanged; \texttt{accounting="hlw"} uses the effective parameters \eqref{eq:prelim-eff} computed by \texttt{hlw\_effective}. The latter evaluates the two integrals of \eqref{eq:prelim-eff} in calendar time by the substitution used in the derivation of Proposition~\ref{prop:prelim-novol}, which removes the $w^{-2/3}$ singularity of \eqref{eq:prelim-nutilde-linear} so that Simpson quadrature suffices.

\paragraph{Tests and what they reproduce.}
The test file runs under \texttt{pytest} or as a plain script. It checks: that the bootstrap returns $P(0,n)=(1+r)^{-n}$ on a flat annual curve and reprices its inputs to $10^{-14}$; that log-linear interpolation keeps a flat curve flat between and beyond the pillars; that $\tau(T+\Delta)=T+\Delta/3$ to $10^{-12}$ for two $(T,\Delta)$ pairs and that $\tau(t)=T+\tfrac\Delta3(1-g(t)^3)$ inside the window; that the at-the-money Bachelier price equals $\sigma\sqrt t/\sqrt{2\pi}$ to $10^{-16}$ and the deep-in and deep-out-of-the-money limits are the intrinsic value and zero; and that every figure of the worked example of Section~\ref{subsec:prelim-worked} is reproduced to the five significant figures printed there, from the clock $\tau^*$ through the at-the-money and $\mp50$ basis-point volatilities at the clock and quoted at $T+\Delta$ to the caplet price and the $7.1\%$ reduction against constant-parameter SABR without decay. The two vol-of-vol accountings are asserted against the same section: the differences from Lemma~\ref{lem:prelim-tc} at $\mp50$ and at the money are tenths of a basis point on the total-variance accounting and below $0.003$ basis points on the effective-parameter accounting, the values the double-precision run gives being $[+0.418,\,+0.125,\,+0.147]$ and $[+0.0023,\,-0.0000,\,-0.0023]$ basis points. The test also asserts the closed forms $\nu_{\mathrm{eff}}^2/\nu^2=1+\tfrac{2}{567}\Delta^3/(\tau^{*3}/3)$ and $\rho_{\mathrm{eff}}\nu_{\mathrm{eff}}/(\rho\nu)=(T\tau^*-T^2/2+\Delta^2/15)/(\tau^{*2}/2)$ to $10^{-9}$, which fixes the quadrature as well as the formulas.

\datanote{The test asserts the closed forms quoted in the gap box after Proposition~\ref{prop:prelim-nueff}, and the effective-parameter figures above are the values the code returns for them. Run: \texttt{cd thesis/code \&\& python3 tests/test\_rates.py} (or \texttt{python3 -m pytest tests/}); six tests, all passing on 9~September 2026 with Python~3.14 and numpy~2.4 and with Python~3.12 under pytest.}

\subsection{The \texttt{certify} package}\label{subsec:appC-certify}

The certificate of Chapter~\ref{ch:certification} is implemented as a second package, \texttt{thesis/code/} \texttt{certify/}: one module, a test file \texttt{certify/tests/test\_certify.py} and a figure script \texttt{certify/} \texttt{fig\_example.py}. It depends on \texttt{numpy} and, for the Beta and normal quantile functions, on \texttt{scipy}; because no system Python used in this work carries \texttt{scipy}, the run command is \texttt{cd thesis/code \&\& uv run --with numpy --with scipy python certify/tests/test\_certify.py} (seven tests, about fifteen seconds). The package regenerates Table~\ref{tab:example} and the series of the block-length sweep figure of Chapter~\ref{ch:sharpness} of Chapter~\ref{ch:certification}.

\paragraph{Blocking.}
\texttt{blocks} lays out the record of Section~\ref{sec:cert-scores} in the pattern of Algorithm~\ref{alg:prelim-blocking}: $n=\lfloor(N-s_0)/(2b)\rfloor$ blocks of length $b$ separated by gaps of length $b$, with the trailing gap held inside the record, so that with no fitting prefix $n=\lfloor N/(2b)\rfloor$. \texttt{deployment\_start} returns the first index of the gap-separated deployment block, and \texttt{block\_scores} sums the signed per-period hedging error over each block, which is the score of Definition~\ref{def:score}, for a single path or an array of paths. The block count is the one Chapter~\ref{ch:certification} uses ($250+240+10=500$ at $N=500$, $b=10$); a count that allows a block to open inside the trailing gap gives thirteen blocks at $b=20$ against the twelve of Table~\ref{tab:example}, and is wrong for that reason.

\paragraph{Threshold and certificate.}
\texttt{conformal\_k} and \texttt{conformal\_threshold} are $k=\lceil(1-\alpha)(n+1)\rceil$ and $\hat q=S_{(k)}$ of Definition~\ref{def:prelim-split}, returning $+\infty$ when $k>n$. \texttt{certificate} returns both forms of Theorem~\ref{thm:cov}: the marginal level $1-\alpha-\dK-(n+1)\beta(b)$ and the calibration-conditional level $b_{n,k}(\delta)-\dK-\beta(b)/\delta'$ with its confidence $1-\delta-n\beta(b)$. Its parts are exposed separately: \texttt{beta\_quantile} is the exact $\Beta(k,n+1-k)$ quantile, \texttt{dkw\_level} the weaker form of Corollary~\ref{cor:prelim-beta-bound}, \texttt{coupling\_cost} the $(n+1)\beta(b)$ of Lemma~\ref{lem:firstblock}. Two invariants of the method are enforced by the interface rather than left to documentation. The mixing coefficient enters as a function argument, \texttt{beta\_fn}, and is never estimated from the sample, because Proposition~\ref{prop:noest} shows that it cannot be. The drift term $\dK$ likewise enters as an argument, either a stress value or the plug-in bound below, because the certificate never sees a deployment score. Read as a sequence of calls, the function is Algorithm~\ref{alg:cert}.

\paragraph{The drift term.}
\texttt{ecdf\_gap} and \texttt{dK\_hat} give the two-sample Kolmogorov distance on the pooled grid; \texttt{W1\_hat} is the empirical $\Wone$ computed as the area between the two distribution functions, which is Lemma~\ref{thm:prelim-w1-identity} in code; \texttt{sqrt\_drift\_bound} is the bound $\sqrt{2L\Wone}$ of Proposition~\ref{prop:w1}. The linear form $L\Wone$ is not implemented, and the test exhibits the pair $\delta_0$ against $U[0,1/L]$ on which it fails.

\paragraph{The blocking curve.}
\texttt{deficit\_curve} evaluates $f(b)=\sqrt{b/N}+Nb^{-(r+1)}$, the expected calibration-conditional shortfall bound of Theorem~\ref{thm:cov}(ii) over the mixing class $\Bclass{r}$, which Chapter~\ref{ch:sharpness} minimises. \texttt{optimal\_block\_length} returns the exact minimiser $b^*=(2(r+1))^{2/(2r+3)}N^{3/(2r+3)}$ and \texttt{optimal\_deficit} its value $f(b^*)=\frac{2r+3}{2(r+1)}(2(r+1))^{1/(2r+3)}N^{-r/(2r+3)}$. The bare balance $N^{3/(2r+3)}$, at which the two terms are equal, is kept under the separate name \texttt{balance\_block\_length} and is not the minimiser: at $N=500$, $r=2$ the minimiser is $23.93$ and the balance $14.35$, a factor $1.67$. No function in the package returns the balance under a name that could be read as optimal. \texttt{expected\_shortfall\_bound} evaluates the same object with the actual constants, $1/(2\sqrt{n+2})+(n+1)\beta(b)$, for a supplied $\beta$; for the AR(1) of Chapter~\ref{ch:certification} with $\varphi=\tfrac12$ it is smallest at $b=11$ ($0.1077$, against $0.1089$ at $b=10$), which is a property of a geometrically mixing process rather than an instance of $b^*$, since geometric mixing lies in every $\Bclass{r}$ and $b^*$ depends on $r$.

\paragraph{Comparison, not theorem.}
\texttt{stationary\_bootstrap\_quantiles} is the stationary bootstrap of \citet{politis1994} applied to the $(1-\alpha)$ quantile of the overlapping block sums. It is an empirical comparator and carries no coverage statement. On the AR(1) example without drift ($200$ replications, $200$ bootstrap resamples) the conformal threshold covers the deployment score with frequency $0.935$, the median bootstrap quantile with $0.860$ and its $95\%$ upper envelope with $0.910$; the mean thresholds are $7.70$, $6.00$ and $7.61$ against a true $P$-quantile of $6.54$. The bootstrap estimates a quantile of the calibration law and under-covers a future block, which is why the thesis states no theorem about it.

\paragraph{The AR(1) example.}
\texttt{ar1\_beta\_bound} takes $\varphi$ and returns the map $k\mapsto\tfrac12\sqrt{-\log(1-\varphi^{2k})}$, the first inequality of Lemma~\ref{prop:prelim-ar1} and the bound Table~\ref{tab:example} uses, in the form \texttt{certificate} expects for \texttt{beta\_fn}; \texttt{ar1\_beta\_exact} evaluates $\beta(k)$ by quadrature, giving $3.108\times10^{-4}$ at $\varphi=\tfrac12$, $k=10$, so the bound is a factor $1.57$ loose. \texttt{ar1\_block\_var} and \texttt{ar1\_block\_sd} return the block variance $v_b$, \texttt{ar1\_drift\_dK} the drift distance $\dK=2\Phi\big(b\mu/(2\sqrt{v_b})\big)-1$, and \texttt{simulate\_ar1} draws stationary paths from a supplied innovation standard deviation, to which the caller passes $\sqrt{1-\varphi^2}$ so that the stationary variance is one; the normalisation is stated with every figure, because it changes every threshold while leaving every coverage number unchanged.

\paragraph{Tests and what they reproduce.}
The seven tests are: the blocking layout and the block count above; that \texttt{ar1\_beta\_bound} dominates \texttt{ar1\_beta\_exact} on a grid of $\varphi$ and $k$; \texttt{test\_table\_3\_1}, which reproduces every cell of Table~\ref{tab:example} at $b\in\{5,10,20\}$ with $\mu=0.05$; a Monte Carlo of the whole procedure at $b=10$ ($10^4$ replications, seed $20260909$) giving realised coverage $0.9113$ with standard error $0.0028$, mean $F_Q(\hat q)=0.9096$ against the coupled value $0.9089$, and $\Prob(F_Q(\hat q)\ge0.7750)=0.9710$ against the certified confidence $0.9378$; the drift bound of Proposition~\ref{prop:w1} on a Gaussian pair, where $\dK=0.1974$, $\hat\dK=0.2005$, $\hat\Wone=2.952$ and $\sqrt{2L\hat\Wone}=0.632$; the minimiser and the balance above; and the bootstrap comparison. \texttt{fig\_example.py} sweeps $b=2,\dots,50$ at $\varphi\in\{0.5,0.8\}$ with $10^4$ replications per point and the same seed and writes \texttt{certify/fig\_example.csv}, the series of the block-length sweep figure of Chapter~\ref{ch:sharpness}: breach frequencies $0.109$, $0.093$, $0.098$ at $b=5,10,20$ for $\varphi=\tfrac12$ and $0.112$, $0.089$, $0.091$ for $\varphi=0.8$. Each lies below the certified breach bound $\alpha+(n+1)\beta(b)+(1-k/(n+1))+\dK$, whose excess over $\alpha$ is $0.925$, $0.129$ and $0.130$ at $\varphi=\tfrac12$ and $8.72$, $1.50$ and $0.186$ at $\varphi=0.8$; almost all the slack is in the coupling term $(n+1)\beta(b)$, which is why Chapter~\ref{ch:sharpness} minimises over $b$ rather than reporting a single block length.

\datanote{Table~\ref{tab:example} and the block-length sweep figure of Chapter~\ref{ch:sharpness} of Chapter~\ref{ch:certification} are regenerated by this package, not computed by hand: \texttt{test\_table\_3\_1} and \texttt{fig\_example.py} at seed $20260909$. The regenerated cells are $\beta(b)=1.563\times10^{-2}$, $4.883\times10^{-4}$, $4.768\times10^{-7}$; $(n+1)\beta(b)=0.7971$, $0.0127$, $6.2\times10^{-6}$; $\dK=0.0299$, $0.0391$, $0.0533$; marginal level $0.0730$, $0.8482$, $0.8467$; $b_{n,k}(0.05)=0.8262$, $0.8239$, $0.7791$; $\beta(b)/\delta'=0.3126$, $0.0098$, $9.5\times10^{-6}$; conditional level $0.4837$, $0.7750$, $0.7258$; confidence $0.1686$, $0.9378$, $0.9500$; and $\E F_Q(\hat q^*)=0.90889$ by quadrature against the $\Beta(24,2)$ density. Every one agrees with the digits printed in Chapter~\ref{ch:certification}, so no hand-arithmetic discrepancy remains to be recorded. The three valid bounds on $\beta(10)$ at $\varphi=\tfrac12$ that appear in the thesis, $4.883\times10^{-4}$ from Lemma~\ref{prop:prelim-ar1}, $5.638\times10^{-4}$ from the $k$-uniform form $|\varphi|^k/(2\sqrt{1-\varphi^2})$ and $6.905\times10^{-4}$ from Proposition~\ref{prop:prelim-ar1}, are not corrected into one another; the code implements the first because that is the one the table uses.}

\gap{Three modules of the reproducibility package are not written. \texttt{strip} (caplet-volatility stripping from cap and floor quotes), \texttt{calibrate} (SABR calibration to a quoted smile with the parameter count of Chapter~\ref{ch:identifiability}) and \texttt{scores} (realised hedging-error scores of a fixed policy) are scheduled for programme months 5--9 and 12--17; nothing in this appendix or in Chapters~\ref{ch:usd} and~\ref{ch:zar} may be read as describing code that exists. Within \texttt{certify}, three limitations are recorded: $\dK$ is an input and is supplied by Chapter~\ref{ch:usd} from observed deployment scores with an assumed density bound $L$; the dependent-data rate for $\hat\Wone$ is not implemented, only the estimator; and \texttt{blocks} discards the tail of the record shorter than $2b$ rather than fitting one further block.}

\figplan{Module diagram, three columns. Left column, implemented: \texttt{rates} (\texttt{curve}, \texttt{clock}, \texttt{sabr}) and \texttt{certify} (blocking, threshold, certificate, drift, blocking curve, AR(1) helpers, bootstrap comparator). Middle column, the thesis objects each module computes, with arrows from the left: bootstrap and forward compounded rate, clock and time-changed caplet, Hagan volatility (Chapter~\ref{ch:prelim}); score, threshold, both forms of the certificate, drift bound, optimal block length (Chapters~\ref{ch:certification} and~\ref{ch:sharpness}). Right column, dashed boxes for the modules not yet written, \texttt{strip}, \texttt{calibrate} and \texttt{scores}, with dashed arrows to the caplet-volatility stripping, SABR calibration and hedging-error scores of Chapters~\ref{ch:usd} and~\ref{ch:zar}, annotated with the programme months in which they are scheduled.}{Code modules and the thesis objects they compute. Solid boxes are implemented and tested at the date of this appendix; dashed boxes are scheduled. The figure records which results of the thesis are reproducible from the package as it stands and which depend on code still to be written. Source: the package layout of Section~\ref{sec:appC-code} and the protocol of Section~\ref{sec:appC-protocol}.\label{fig:appC-modules}}


\section{Reproducibility protocol}\label{sec:appC-protocol}

The protocol fixes nine things: where the package lives, what environment runs it, which seeds it uses, what the manifest records, how an examiner verifies results whose inputs cannot be shipped, how quoted and derived quotes are separated, which outputs are expected and to what tolerance, how sponsor data are governed, and what tool assistance was used.

\paragraph{1. Package and paths.} The package is one directory, \texttt{Code/}, at the root of the thesis repository, holding \texttt{rates/}, \texttt{certify/}, \texttt{tests/} and \texttt{data/}; the paths quoted above as \texttt{thesis/code/\ldots} denote it, and commands are written relative to it. Nothing outside \texttt{Code/} is needed to reproduce any number in Chapters~\ref{ch:prelim} to~\ref{ch:sharpness}.

\gap{The archive of record for the package, its persistent identifier and the licence under which it is released are to be fixed at first submission, together with the statement in Chapter~\ref{ch:conclusions} of what the archive contains; they are not yet minted.}

\paragraph{2. Environment.} One environment runs everything: Python~3.12 or later, \texttt{numpy}, and \texttt{scipy} for the quantile functions \texttt{certify} needs. The \texttt{rates} package is \texttt{numpy}-only by design, taking the normal distribution function from \texttt{math.erf}, so Chapter~\ref{ch:prelim} reproduces without \texttt{scipy}. The two commands are
\begin{quote}\small
\texttt{cd Code \&\& python3 tests/test\_rates.py}\\
\texttt{cd Code \&\& uv run --with numpy --with scipy python certify/tests/test\_certify.py}
\end{quote}
the second in that form because no system Python used in this work carries \texttt{scipy}. Verified on 9~September 2026: Python~3.14.6 with numpy~2.4.6, and Python~3.12 under \texttt{pytest}, for the six \texttt{rates} tests; \texttt{scipy} through an ephemeral \texttt{uv} environment for the seven \texttt{certify} tests.

\gap{The version pin itself, a \texttt{requirements.txt} with exact versions and a lock file, is not yet in the package; until it is, the versions recorded in the paragraph above are the only pin.}

\paragraph{3. Seeds.} Every computed result in the thesis comes from a run in Table~\ref{tab:appC-seeds}. The seed is set once per script and the generator passed explicitly, so a test run alone gives the numbers it gives inside the suite.

\begin{table}[htbp]\centering\small
\caption{Random and deterministic runs behind the computed results of the thesis, with script, seed, replication count and approximate runtime on a laptop. Deterministic entries have no seed.}\label{tab:appC-seeds}
\begin{tabular}{@{}p{4.3cm}p{4.6cm}p{1.9cm}rr@{}}
\toprule
Result & Script & Seed & Reps & Time \\
\midrule
Worked example, Section~\ref{subsec:prelim-worked} & \texttt{tests/test\_rates.py} & --- & --- & $<1$\,s \\
Table~\ref{tab:example} & \texttt{certify/tests/test\_certify.py} & --- & --- & $<1$\,s \\
Coverage Monte Carlo, Section~\ref{sec:cert-example} & \texttt{certify/tests/test\_certify.py} & 20260909 & $10^4$ & $10$\,s \\
Drift-bound check, Proposition~\ref{prop:w1} & \texttt{certify/tests/test\_certify.py} & 20260910 & 4400 & $<1$\,s \\
Bootstrap comparison & \texttt{certify/tests/test\_certify.py} & 20260911 & 200 & $3$\,s \\
the block-length sweep figure of Chapter~\ref{ch:sharpness} series & \texttt{certify/fig\_example.py} & 20260909 & $10^4$/point & $6$\,s \\
\bottomrule
\end{tabular}
\end{table}

\paragraph{4. Data manifest.} One row per dataset in \texttt{Code/data/manifest.csv}: identifier (ticker, RIC, endpoint or landing page); for every Reserve Bank export, the selector settings, that is the date range requested, the earliest date returned and the download timestamp, since those pages serve no static file; local file name; SHA-256 checksum; access class; and in-package flag, one of \emph{yes}, \emph{script and hash}, \emph{hash only}, \emph{no}. The registers of Sections~\ref{sec:appC-usd} and~\ref{sec:appC-zar} fix the class and flag of every row in advance; the manifest adds checksums as files appear.

\paragraph{5. Verification without the shipped data.} An examiner with terminal access re-exports each class~T row by its identifier and compares the checksum. A mismatch is expected for composite volatility surfaces, which vendors revise, so the manifest carries the export date and the vendor's revision policy beside the hash and a mismatch is reported, not resolved silently. An examiner without terminal access reproduces the whole of Chapters~\ref{ch:prelim} to~\ref{ch:sharpness}, Table~\ref{tab:example} and the block-length sweep figure of Chapter~\ref{ch:sharpness} included, together with the class~F and class~C rows of both registers.

\paragraph{6. Quoted against derived.} Every volatility point carries a flag, \emph{quoted}, \emph{derived} or \emph{interpolated}, from the vendor field where one exists and otherwise by the rule of Section~\ref{appC:usd-libor-vol}, under which strikes beyond ATM\,$\pm100$~bp count as unquoted. No result aggregates quoted and derived points without the flag.

\paragraph{7. Expected outputs and tolerances.} The tests assert printed figures, not merely that the code runs: the Chapter~\ref{ch:prelim} worked example to five significant figures, every cell of Table~\ref{tab:example} to its four decimals, the effective-parameter closed forms to $10^{-9}$, the clock identity to $10^{-12}$. Monte Carlo results are asserted within three standard errors, so the coverage $0.9113$ reproduces to within $0.0084$, and the figure series are compared with the committed \texttt{fig\_example.csv}.

\paragraph{8. Sponsor data.} Governance is stated once, in Section~\ref{sec:appC-vols}. In the manifest the sponsor rows carry the flag \emph{no} and no checksum is published.

\paragraph{9. Tool assistance.} Generative artificial-intelligence assistance was used in this work and is declared as the University's rules require; the categories are code development, presentation and editing, and citation formatting. Every number so produced is verified by a run in Table~\ref{tab:appC-seeds} or by independent arithmetic recorded in the data note of the chapter that prints it, every citation is checked against the cited document, and no mathematical statement is admitted on the strength of tool output alone, which is why unproved statements appear as proof routes with an explicit list of what remains to be established.

\gap{The tool names, versions and dates of use, and the mapping to the categories of the Faculty declaration form, are to be listed here at first submission; the record is kept as the work proceeds and is not yet reproduced.}

\subsection{Terminal session checklist}\label{subsec:appC-checklist}

The identifiers and quantities marked to be confirmed in Sections~\ref{sec:appC-usd} and~\ref{sec:appC-zar} are closed in one terminal session. Its required returns are listed here so that the count can be driven to zero and the registers then cited as verified: the first calendar month in which the volatility cube marks the SOFR cap surface as quoted rather than derived; the cap and swaption volatility ticker roots for 3M~LIBOR and for compounded SOFR, with their strike and expiry encoding; the SOFR OIS strip identifiers and the curve number carrying them; the history depth and tenor points of the LIBOR--SOFR basis-swap curve; the 1W and 2M USD fallback spread values; the front SOFR futures identifier; the LIBOR fixing identifier and its history depth; and, for South Africa, the curve identifiers for the live ZARONIA OIS and JIBAR swap curves, the tenor columns of the exchange's daily files, and the earliest date the Reserve Bank selector returns for ZARONIA and for JIBAR. Each return enters the manifest with the session date and replaces one to-be-confirmed entry in the register.
